\exercisegroup

\begin{enumerate}
\def\labelenumi{\arabic{enumi})}
\item
  证明：Hausdorff 桶型空间的完备化是桶型的。
\item
  设 E 是 R 或 C 上的一个向量空间。证明：赋以 E 上最细局部凸拓扑（II, 第 25 页）的 E 是桶型的。由此给出不可度量化且不是 Baire 空间的桶型空间的例子。
\item
  设 \(E\) 是一个 Hausdorff 局部凸空间，具有可数无穷基 \((a_{n})\) 。
\end{enumerate}

\begin{enumerate}
\def\labelenumi{\alph{enumi})}
\item
  证明：E 具有一个可数的、拓扑无关的基 \((e_n)\) （利用 E 中每条直线都有拓扑补这一事实，用归纳法定义诸 \(e_n\) ）。
\item
  证明：要使 E 是桶型的，必须且只需 E 的拓扑 T 与 E 上的最细局部凸拓扑相同（注意每个序列 \((\lambda_{n}e_{n})\) 的均衡凸包在 E 中闭）。特别地，若 T 可度量化，则 E 不是桶型的（cf.~习题 2）。
\end{enumerate}

\begin{enumerate}
\def\labelenumi{\arabic{enumi})}
\setcounter{enumi}{3}
\item
  设 E 是一个 Banach 空间，其中存在一个无穷的代数无关序列 \((a_{n})\) ，它在 E 中是完全的（例如空间 \(\ell^{1}(\mathbf{N})\) （I, 第 4 页））。设 B 是 E 的一个包含诸 \(a_{n}\) 的基；我们知道（II, 第 80 页, 习题 24）B 不可数。设 \((e_{n})\) 是 B 中互异且异于诸 \(a_{n}\) 的元素组成的序列，并设 C 是诸 \(e_{n}\) 组成的集在 B 中的余集。设 \(F_{n}\) 是 E 中由 C 与诸 \(e_{k}\) （指标 \(k \leqslant n\) ）生成的向量子空间；E 是诸 \(F_{n}\) 的并。设 S 是 E 中的单位球；证明存在一个指标 n，使得 \(S \cap F_{n}\) 不是第一纲集。由此推出：对这个 n 值，\(F_{n}\) 是一个可度量化的、非完备的 Baire 空间。
\item
  给出一个局部凸空间的例子，它是完备的 Hausdorff Baire 空间，但不可度量化（cf.~GT, IX, § 5, 习题 16）。
\item
  称局部凸空间 \(\mathbf{E}\) 是相对有界的，如果 \(\mathbf{E}\) 中存在一个有界的桶。
\end{enumerate}

\begin{enumerate}
\def\labelenumi{\alph{enumi})}
\item
  要使 E 相对有界，必须且只需 E 的拓扑比某个由半范数定义的拓扑粗。此时存在 E 的典范有界结构的一个由桶组成的基。
\item
  要使 E 既是有界型的又是相对有界的，必须且只需 E 的拓扑是 E 上一族赋范空间拓扑的下确界（cf.~III, 第 40 页, 习题 1）。此外，要使 E 的典范有界结构还具有可数基，必须且只需 E 的拓扑是赋范空间拓扑的一个序列的下确界。
\end{enumerate}

\begin{enumerate}
\def\labelenumi{\arabic{enumi})}
\setcounter{enumi}{6}
\item
  称局部凸空间 E 是次桶型的，如果 E 的每个吸有界集的（III, 第 39 页, 习题 11）桶都是 E 中 0 的邻域。每个有界型空间都是次桶型的；每个桶型空间都是次桶型的。证明：Hausdorff 次桶型空间的完备化是桶型的（利用如下事实：在 Hausdorff 局部凸空间 E 中，每个桶都吸收 E 的每个凸的、均衡的、有界的且半完备的子集）。
\item
  设 \((\mathrm{E}_{\iota})_{\iota\in I}\) 是一族次桶型空间，且对每个 \(\iota\in I\) ，设 \(f_{\iota}\) 是从 \(\mathrm{E}_{\iota}\) 到一个向量空间 E 中的线性映射。证明：赋以使诸 \(f_{\iota}\) 连续的最细局部凸拓扑的空间 E 是次桶型的。特别地，次桶型空间的每个商空间都是次桶型的；次桶型空间的每个拓扑直和都是次桶型的。
\item
  设 I 是一个不可数无穷集；在直和向量空间 \(E = R^{(I)}\) 上，一方面考虑最细局部凸拓扑 T，另一方面考虑 I, 第 24 页, 习题 14 中定义的拓扑 \(T_{0}\) ，它是局部凸的；我们知道 T 与 \(T_{0}\) 不同（II, 第 75 页, 习题 11），且赋以 \(T_{0}\) 的 E 是完备的（GT, III, § 3, 习题 10）。证明：E 中的有界集对 T 与 \(T_{0}\) 是相同的（III, 第 39 页, 习题 10），且赋以 \(T_{0}\) 的 E 不是桶型的（注意由所有 \(x = (\xi_{\mathfrak{u}}) \in \mathbf{E}\) （满足 \(\sum_{\mathfrak{u} \in I} |\xi_{\mathfrak{u}}| \leqslant 1\) ）组成的集 T 是一个桶，并利用 II, 第 75 页, 习题 11）。
\item
  证明：若一个次桶型空间中每个闭的凸均衡有界子集都是半完备的，则它是一个桶型空间。
\item
  设 \(\mathbf{E}\) 是一个次桶型空间，\(\mathbf{F}\) 是一个局部凸空间。证明：\(\mathcal{L}(\mathbf{E};\mathbf{F})\) 中对有界收敛拓扑有界的每个子集都等度连续。
\end{enumerate}

¶ 12) a) 设 E 是一个局部凸空间，\((A_n)\) 是 E 中凸均衡集的递增序列，使得 \(A = \bigcup_{n} A_n\) 是吸收的。设 \((W_n)\) 是凸均衡的

0 的邻域组成的递减序列；那么诸 \(\mathrm{W}_n\cap \mathrm{A}_n\) 的均衡凸包 V 是吸收的；若 E 是桶型的，则 \(\overline{\mathbf{V}}\) 是 0 的一个邻域。

\begin{enumerate}
\def\labelenumi{\alph{enumi})}
\setcounter{enumi}{1}
\item
  设 \(\mathfrak{F}\) 是 E 上的一个滤子；假设对每个 \(n\) ，存在一个集 \(\mathrm{M}_n \in \mathfrak{F}\) ，使得 \((\mathrm{M}_n + \mathrm{W}_n) \cap \mathrm{A}_{2n} = \emptyset\) 。设 \(\mathrm{V}_n\) 是诸 \(\mathrm{W}_k \cap \mathrm{A}_k\) （\(k \leqslant n - 1\) ）与 \(\mathrm{W}_n\) 的均衡凸包，使得 \(\mathrm{V}_n\) 是 0 的一个邻域，且有 \(\frac{1}{2}\overline{\mathrm{V}} \subset \mathrm{V}_n\) （对一切 \(n\)）。证明 \((\mathrm{M}_n + \mathrm{V}_n) \cap \mathrm{A}_n = \emptyset\) （对一切 \(n\) ）。
\item
  由 a) 与 b) 推出：若 E 是桶型的且 \(\mathfrak{F}\) 是 E 上的一个 Cauchy 滤子，则存在整数 N，使得对一切 \(\mathbf{M} \in \mathfrak{F}\) 以及 E 中 0 的每个邻域 W，\(\mathbf{M} + \mathbf{W}\) 都与 \(\mathrm{A}_{\mathbf{N}}\) 相交。（用反证法；采用 b) 的记号，考虑一个集合 \(\mathbf{M} \in \mathfrak{F}\) （具有小阶 \(\frac{1}{2}\overline{\mathrm{V}}\)）。）
\end{enumerate}

¶ 13) a) 设 E 是一个桶型空间，\((C_n)\) 是凸的、均衡的集的递增序列，使得 \(E = \bigcup_{n} C_n\) 。设 U 是一个凸的、均衡的且吸收的集，使得对每个 n，

\(U \cap C_{n}\) 在 \(C_{n}\) 中闭。证明 U 是 E 中 0 的一个邻域。（证明 \(\overline{U} \subset 2U\) ，为此考虑 U 上的一个滤子 \(\mathfrak{F}\) （收敛于一点 \(x \in E\) ）并应用习题 12, c)。）

\begin{enumerate}
\def\labelenumi{\alph{enumi})}
\setcounter{enumi}{1}
\tightlist
\item
  设 \(\mathrm{E}\) 是一个桶型空间，\((\mathrm{E}_n)\) 是 \(\mathbf{E}\) 的子空间的递增序列，使得 \(\mathrm{E} = \bigcup_{n} \mathrm{E}_n\) 。
\end{enumerate}

证明：若 \(\mathbf{U}\) 是 \(\mathbf{E}\) 的一个子集，使得 \(\mathbf{U} \cap \mathbf{E}_n\) 是 \(\mathbf{E}_n\) 中的一个桶（对每个 \(n\) ），则 \(\mathbf{U}\) 是 \(\mathbf{E}\) 中 0 的一个邻域。特别地，\(\mathbf{E}\) 是序列 \((\mathbf{E}_n)\) 的严格归纳极限（II, 第 33 页）。

¶ 14) a) 设 E 是一个 Hausdorff 局部凸空间，L 是 E 的一个有限余维数的子空间，T 是 L 中的一个桶。证明：存在 E 中的一个桶 T'，使得 T' ∩ L = T（证明可以取 T' 为 T 在 E 中的闭包 T̅ 与一个有限维紧凸集之和）。

\begin{enumerate}
\def\labelenumi{\alph{enumi})}
\setcounter{enumi}{1}
\tightlist
\item
  设 E 是一个桶型空间，L 是 E 的一个子空间，它有一个具有可数基的补。证明 L 是桶型的（利用 \(a\) ) 与习题 13, b)）。
\end{enumerate}

* c) 设 E 是一个 Hausdorff 局部凸空间；它的完备化 \(\hat{\mathbf{E}}\) 可等同于一个桶型空间 F 的闭子空间，F 是一族 Fréchet 空间的乘积（II, 第 5 页, 命题 3 与 IV, 第 14 页, 推论）。设 \((e_{\alpha})_{\alpha \in \mathbf{A}}\) 是子空间 E 在 F 中的一个补的基，并设 \(\mathrm{H}_{\alpha}\) 是 F 中由 E 与诸 \(e_{\beta}\) （指标 \(\beta \neq \alpha\) ）生成的超平面；由 \(b\) )，\(\mathrm{H}_{\alpha}\) 是桶型空间。对每个 \(x \in \mathrm{E}\) ，令 \(u(x)\) 为桶型空间 \(\mathrm{G} = \prod_{\alpha \in \mathbf{A}} \mathrm{H}_{\alpha}\) （

\(F^{A}\) 的子空间）中所有坐标都等于 x 的点；u 是从 E 到子空间 \(\Delta \cap G\) 上的同构，其中 \(\Delta\) 是 \(F^{A}\) 中的对角线。证明 \(u(E) = \Delta \cap G\) 在 G 中闭，从而每个 Hausdorff 局部凸空间都同构于一个 Hausdorff 桶型空间的闭子空间。*

\begin{enumerate}
\def\labelenumi{\arabic{enumi})}
\setcounter{enumi}{14}
\tightlist
\item
  设 E 是一个 Hausdorff 桶型（相应地，次桶型）空间，\(\hat{E}\) 是它的完备化。证明：\(\hat{E}\) 的每个包含 E 的子空间 F 都是桶型的（相应地，次桶型的）（cf.~III, 第 24 页, 推论 与 IV, 第 52 页, 习题 1）。
\end{enumerate}

¶ * 16) 设 \((E_{\iota})_{\iota\in I}\) 是不可数族 Hausdorff 桶型空间，其中没有一个是点 0，并设 \(E = \prod_{\iota \in I} E_{\iota}\) ；那么 \(E\) 是桶型的（IV, 第 14 页, 推论）。设 \(G\) 是 E 中由所有点 \((x_{\iota})\) （除可数多个指标外均有 \(x_{\iota}=0\) ）组成的子空间。G 中每条在 E 中收敛的点列的极限都属于 G，但 G 在 E 中稠密。

\begin{enumerate}
\def\labelenumi{\alph{enumi})}
\tightlist
\item
  证明：每个子集 \(\mathbf{M}\) （\(\mathrm{G}' = \mathrm{E}'\) 的）若对 \(\sigma(\mathrm{E}', \mathrm{G})\) 有界，都含于一个有限乘积 \(\prod_{\iota \in \mathrm{H}} \mathrm{E}_{\iota}'\) （IV, 第 12 页, 命题 13）之中，其中 \(\mathrm{H}\) 是 I 的有限子集；因而 \(\mathbf{M}\)
\end{enumerate}

对 \(\sigma(E', E)\) 有界。由此推出 \(G\) 是桶型的。

\begin{enumerate}
\def\labelenumi{\alph{enumi})}
\setcounter{enumi}{1}
\tightlist
\item
  设 F 是 E 的一个子空间，使得 \(G \subset F \subset E\) ，且 G 是 F 中的一个（处处稠密的）超平面；F 是桶型的（习题 15）。证明 F 不是有界型的。（用反证法；若 F 中存在一个凸的、均衡的且有界的集 A，使得 G 是赋范空间 \(F_{A}\) （III, 第 7 页）中的处处稠密超平面，则将存在 G 的一个点列收敛于 F 中不属于 G 的一点）（cf.~IV, 第 52 页, 习题 2）。*
\end{enumerate}

\begin{enumerate}
\def\labelenumi{\arabic{enumi})}
\setcounter{enumi}{16}
\tightlist
\item
  \begin{enumerate}
  \def\labelenumii{\alph{enumii})}
  \tightlist
  \item
    设 E 是一个 Hausdorff 局部凸空间，L 是 E 的一个有限余维数的向量子空间，T 是 L 中一个吸有界集的桶。证明：存在 E 中一个吸有界集的桶 \(T'\) ，使得 \(T' \cap L = T\) 。（化归到 L 是 E 中超平面的情形。设 \(E_{0}\) 是与 E 相伴的有界型空间（III, 第 40 页, 习题 1），\(L_{0}\) 是赋以 \(E_{0}\) 的拓扑所诱导拓扑的超平面 L；注意 T 是 \(L_{0}\) 中 0 的一个邻域，并考虑以下两种情形：\(L_{0}\) 在 \(E_{0}\) 中稠密，或在 \(E_{0}\) 中闭；证明对 \(T'\) 可取 T 在 E 中的闭包 \(\overline{T}\) ，或 T 与一个 1 维紧凸集之和）。
  \end{enumerate}
\end{enumerate}

\begin{enumerate}
\def\labelenumi{\alph{enumi})}
\setcounter{enumi}{1}
\tightlist
\item
  设 E 是一个次桶型空间，L 是 E 的一个有限余维数的向量子空间。由 a) 推出 L 是次桶型的（cf.~IV, 第 64 页, 习题 11）。
\end{enumerate}

¶ 18) 设 E 是局部凸可度量化子空间的递增序列 ( \(E_{n}\) ) 的严格归纳极限空间（II, 第 33 页），并设 F 是 E 的一个向量子空间，使得 E 的每一点都是 F 的某个点列的极限点。

\begin{enumerate}
\def\labelenumi{\alph{enumi})}
\item
  若 \(\overline{\mathbf{E}}_n\) 是 \(\mathbf{E}_n\) 在 \(\mathbf{E}\) 中的闭包，则 \(\mathbf{E}\) 是序列 \((\overline{\mathbf{E}}_n)\) 的严格归纳极限。设 \(\mathbf{F}_n\) 是 \(\mathbf{F} \cap \overline{\mathbf{E}}_n\) 在 \(\mathbf{E}\) 中的闭包。证明 \(\mathbf{E}\) 是子空间的递增序列 \(\mathbf{F}_n\) 的并。
\item
  假设 E 是桶型的。证明 F 是有界型的。（设 u 是从 F 到一个 Banach 空间 G 中的线性映射，它把 F 的每个有界子集变为 G 的有界子集。证明存在从 E 到 G 中的一个线性映射，它在 F 上的限制等于 u，且在每个 \(F_{n}\) 上的限制连续。最后利用 III, 第 44 页, 习题 13, b)。）
\end{enumerate}

\begin{enumerate}
\def\labelenumi{\arabic{enumi})}
\setcounter{enumi}{18}
\tightlist
\item
  称 Hausdorff 局部凸空间 E 是超有界型的，如果 E 的每个吸收 E 的所有凸的、均衡的、有界的且半完备的子集的凸子集，都是 E 中 0 的邻域。
\end{enumerate}

\begin{enumerate}
\def\labelenumi{\alph{enumi})}
\item
  证明：每个超有界型空间都既是有界型的又是桶型的。
\item
  设 E 是一个 Hausdorff 局部凸空间，使得每个趋于 0 的序列的点所成之集的闭均衡凸包都是半完备的。证明：若 E 是有界型的，则它是超有界型的。特别地，每个有界型的拟完备空间都是超有界型的；每个 Fréchet 空间都是超有界型的。
\item
  设 \((\mathrm{E}_{\alpha})\) 是向量空间 E 的向量子空间的一个有向递增族，使得 E 是诸 \(\mathrm{E}_{\alpha}\) 的并。设 \(\mathcal{T}_{\alpha}\) 是 \(\mathrm{E}_{\alpha}\) 上的一个局部凸拓扑（对每个 \(\alpha\) ），并设 \(\mathcal{T}\) 是使从 \(\mathrm{E}_{\alpha}\) 到 E 中的诸典范单射都连续的最细局部凸拓扑。假设 \(\mathcal{T}\) 是 Hausdorff 的，且对每个 \(\alpha\) ，\(\mathrm{E}_{\alpha}\) 上由 \(\mathcal{T}\) 诱导的拓扑是 \(\mathcal{T}_{\alpha}\) 。证明：若诸空间 \(\mathrm{E}_{\alpha}\) 中的每一个都是超有界型的，则赋以 \(\mathcal{T}\) 的 E 是超有界型的。d) 证明：超有界型空间的每个有限乘积都是超有界型的；由此推出超有界型空间的每个拓扑直和都是超有界型的。
\item
  证明：超有界型空间的无穷序列的乘积空间 \(\mathrm{E} = \prod_{n=0}^{\infty} \mathrm{E}_n\)
\end{enumerate}

是超有界型的。（设 A 是 E 的一个凸子集，它吸收 E 的每个凸的、均衡的、有界的且半完备的子集。证明：若 A 不是 E 中 0 的邻域，则将存在一个序列 \((x_{n})\) （在 \(\mathbb{C}\) A 中），使得 \(x_{n}\) 的前 \(n - 1\) 个坐标为零，但 \(\neq 0\) 。然后注意，这样一个序列的点所成之集的闭均衡凸包等同于点 \(\sum_{n = 0}^{\infty}\lambda_n x_n\) （其中 \(\sum_{n = 0}^{\infty}|\lambda_n| \leqslant 1\) ）所成的集，且

这个包是一个半完备的集。）

\begin{enumerate}
\def\labelenumi{\arabic{enumi})}
\setcounter{enumi}{19}
\item
  证明：要使 Hausdorff 局部凸空间 \(E\) 是超有界型的，必须且只需它是 Banach 空间的一个族的归纳极限。（为看条件必要，考虑凸的、均衡的、有界的且半完备的集 \(B\) （在 \(E\) 中）以及空间 \(E_B\) 。为看条件充分，注意：若 \(E\) 是 Banach 空间的一个族 \(E_\alpha\) 的归纳极限，则可以假设诸 \(E_\alpha\) 是 \(E\) 的（代数意义上的）子空间；若 \(V\) 是 \(E\) 中吸收 \(E\) 的凸的、均衡的、有界的且半完备的子集的一个凸集，证明 \(V\) 吸收每个球 \(B_\alpha\) （\(E_\alpha\) 的）（用反证法）；若 \(V\) 不吸收 \(B_\alpha\) ，则它不吸收一个点列 \((x_n)\) （由 \(B_\alpha\) 的点组成，在 \(E\) 中趋于 0）；然后利用如下事实：在 Banach 空间中，紧集的闭凸包是紧的。）
\item
  证明：若 E 是 Hausdorff 局部凸半完备空间，则与 E 相伴的有界型空间（III, 第 40 页, 习题 1）是超有界型的。
\end{enumerate}

¶ 22) 设 E 是一个满足第一可数公理的无穷维 Banach 空间。

\begin{enumerate}
\def\labelenumi{\alph{enumi})}
\item
  证明：集 \(\mathcal{K}\) （由 E 中使 \(\mathrm{E_A}\) 为无穷维的所有紧凸均衡子集 A 组成）是无穷的，且其基数 \(\leqslant 2^{\mathrm{Card}(\mathbf{N})}\) （GT, IX, § 5, 习题 17）。对每个 \(x_0 \in \mathrm{E}\) 及每个 \(\mathrm{A} \in \mathcal{K}\) ，集 \(x_0 + \mathrm{A}\) 都含有一个基数为 \(2^{\mathrm{Card}(\mathbf{N})}\) 的自由子集（II, 第 80 页, 习题 24, c)）。
\item
  设 E 中 \(x_0 \neq 0\) 。证明：存在一个族 \((y_{\mathrm{A}})_{\mathrm{A} \in \mathcal{K}}\) ，使得 \(x_0\) 与诸 \(y_{\mathrm{A}}\) 构成一个自由族，且有 \(y_{\mathrm{A}} \in x_0 + \mathrm{A}\) （对一切 \(\mathrm{A} \in \mathcal{K}\) ；将 \(\mathcal{K}\) 良序化，并利用 \(a\) ) 作超限归纳论证）。
\item
  设 \(f \in \mathrm{E}^{*}\) 是满足 \(f(x_0) = 1\) 且 \(f(y_{\mathrm{A}}) = 0\) （对一切 \(\mathrm{A} \in \mathcal{K}\)）的线性型，并设 \(\mathrm{H} = f^{-1}(0)\) 。证明：一个子集 \(\mathbf{M}\) （\(\mathrm{H}\) 的）若是凸的、均衡的且半完备的，则必是有限维的（注意否则 \(\mathbf{M}\) 将含有一个无穷维的紧凸均衡集 \(\mathrm{A}\) ；于是 \(y_{\mathrm{A}}\) 将属于 \(\mathrm{H} \cap (x_0 + \mathbf{M})\) ）。
\item
  证明：赋以 E 的拓扑所诱导拓扑的 H 不是超有界型的，尽管它是有界型的且是桶型的（III, 第 44 页, 习题 14）。（利用 c)，证明若 H 是超有界型的，则其拓扑将是最细局部凸拓扑，由此得出矛盾。）
\end{enumerate}

